%=====================================
%   semiclassical.tex
%   Semiclassical Analysis 
%   toshi2019
%=====================================

% -----------------------
% preamble
% -----------------------
% ここから本文 (\begin{document}) までの
% ソースコードに変更を加えた場合は
% 編集者まで連絡してください. 
% Don't change preamble code yourself. 
% If you add something
% (usepackage, newtheorem, newcommand, renewcommand),
% please tell it 
% to the editor of institutional paper of RUMS.

% ------------------------
% documentclass
% ------------------------
\documentclass[leqno, 10pt, a4paper, dvipdfmx]{jsarticle}

% ------------------------
% usepackage
% ------------------------
\usepackage{algorithm}
\usepackage{algorithmic}
\usepackage{amscd}
\usepackage{amsfonts}
\usepackage{amsmath}
\usepackage[psamsfonts]{amssymb}
\usepackage{amsthm}
\usepackage{ascmac}
\usepackage{bm}
\usepackage{color}
\usepackage{enumerate}
\usepackage{fancybox}
\usepackage[stable]{footmisc}
\usepackage{graphicx}
\usepackage{listings}
\usepackage{mathrsfs}
\usepackage{mathtools}
\usepackage{otf}
\usepackage{pifont}
\usepackage{proof}
\usepackage{subfigure}
\usepackage{tikz}
\usepackage{verbatim}
\usepackage[all]{xy}

\usetikzlibrary{cd}
\usetikzlibrary{arrows.meta}


% ================================
% パッケージを追加する場合のスペース 
\usepackage[dvipdfmx]{hyperref}
\usepackage{xcolor}
\definecolor{darkgreen}{rgb}{0,0.45,0} 
\definecolor{darkred}{rgb}{0.75,0,0}
\definecolor{darkblue}{rgb}{0,0,0.6} 
\hypersetup{
    colorlinks=true,
    citecolor=darkgreen,
    linkcolor=darkred,
    urlcolor=darkblue,
}
\usepackage{pxjahyper}

%\usepackage[mathscr]{euscript} % mathscr のフォントを変更

%\usepackage{mmasym}

%=================================


% --------------------------
% theoremstyle
% --------------------------
\theoremstyle{definition}

% --------------------------
% newtheoem
% --------------------------

% 日本語で定理, 命題, 証明などを番号付きで用いるためのコマンドです. 
% If you want to use theorem environment in Japanece, 
% you can use these code. 
% Attention!
% All theorem enivironment numbers depend on 
% only section numbers.
\newtheorem{Axiom}{公理}[section]
\newtheorem{Definition}[Axiom]{定義}
\newtheorem{Theorem}[Axiom]{定理}
\newtheorem{Proposition}[Axiom]{命題}
\newtheorem{Lemma}[Axiom]{補題}
\newtheorem{Corollary}[Axiom]{系}
\newtheorem{Example}[Axiom]{例}
\newtheorem{Claim}[Axiom]{主張}
\newtheorem{Property}[Axiom]{性質}
\newtheorem{Attention}[Axiom]{注意}
\newtheorem{Question}[Axiom]{問}
\newtheorem{Problem}[Axiom]{問題}
\newtheorem{Consideration}[Axiom]{考察}
\newtheorem{Alert}[Axiom]{警告}
\newtheorem{Fact}[Axiom]{事実}


% 日本語で定理, 命題, 証明などを番号なしで用いるためのコマンドです. 
% If you want to use theorem environment with no number in Japanese, You can use these code.
\newtheorem*{Axiom*}{公理}
\newtheorem*{Definition*}{定義}
\newtheorem*{Theorem*}{定理}
\newtheorem*{Proposition*}{命題}
\newtheorem*{Lemma*}{補題}
\newtheorem*{Example*}{例}
\newtheorem*{Corollary*}{系}
\newtheorem*{Claim*}{主張}
\newtheorem*{Property*}{性質}
\newtheorem*{Attention*}{注意}
\newtheorem*{Question*}{問}
\newtheorem*{Problem*}{問題}
\newtheorem*{Consideration*}{考察}
\newtheorem*{Alert*}{警告}
\newtheorem{Fact*}{事実}


% 英語で定理, 命題, 証明などを番号付きで用いるためのコマンドです. 
%　If you want to use theorem environment in English, You can use these code.
%all theorem enivironment number depend on only section number.
\newtheorem{Axiom+}{Axiom}[section]
\newtheorem{Definition+}[Axiom+]{Definition}
\newtheorem{Theorem+}[Axiom+]{Theorem}
\newtheorem{Proposition+}[Axiom+]{Proposition}
\newtheorem{Lemma+}[Axiom+]{Lemma}
\newtheorem{Example+}[Axiom+]{Example}
\newtheorem{Corollary+}[Axiom+]{Corollary}
\newtheorem{Claim+}[Axiom+]{Claim}
\newtheorem{Property+}[Axiom+]{Property}
\newtheorem{Attention+}[Axiom+]{Attention}
\newtheorem{Question+}[Axiom+]{Question}
\newtheorem{Problem+}[Axiom+]{Problem}
\newtheorem{Consideration+}[Axiom+]{Consideration}
\newtheorem{Alert+}{Alert}
\newtheorem{Fact+}[Axiom+]{Fact}
\newtheorem{Remark+}[Axiom+]{Remark}

% ----------------------------
% commmand
% ----------------------------
% 執筆に便利なコマンド集です. 
% コマンドを追加する場合は下のスペースへ. 

% 集合の記号 (黒板文字)
\newcommand{\NN}{\mathbb{N}}
\newcommand{\ZZ}{\mathbb{Z}}
\newcommand{\QQ}{\mathbb{Q}}
\newcommand{\RR}{\mathbb{R}}
\newcommand{\CC}{\mathbb{C}}
\newcommand{\PP}{\mathbb{P}}
\newcommand{\KK}{\mathbb{K}}


% 集合の記号 (太文字)
\newcommand{\nn}{\mathbf{N}}
\newcommand{\zz}{\mathbf{Z}}
\newcommand{\qq}{\mathbf{Q}}
\newcommand{\rr}{\mathbf{R}}
\newcommand{\cc}{\mathbf{C}}
\newcommand{\pp}{\mathbf{P}}
\newcommand{\kk}{\mathbf{k}}

% 特殊な写像の記号
\newcommand{\ev}{\mathop{\mathrm{ev}}\nolimits} % 値写像
\newcommand{\pr}{\mathop{\mathrm{pr}}\nolimits} % 射影

% スクリプト体にするコマンド
%   例えば　{\mcal C} のように用いる
\newcommand{\mcal}{\mathcal}

% 花文字にするコマンド 
%   例えば　{\h C} のように用いる
\newcommand{\h}{\mathscr}

% ヒルベルト空間などの記号
\newcommand{\F}{\mcal{F}}
\newcommand{\X}{\mcal{X}}
\newcommand{\Y}{\mcal{Y}}
\newcommand{\Hil}{\mcal{H}}
\newcommand{\RKHS}{\Hil_{k}}
\newcommand{\Loss}{\mcal{L}_{D}}
\newcommand{\MLsp}{(\X, \Y, D, \Hil, \Loss)}

% 偏微分作用素の記号
\newcommand{\p}{\partial}

% 角カッコの記号 (内積は下にマクロがあります)
\newcommand{\lan}{\langle}
\newcommand{\ran}{\rangle}



% 圏の記号など
\newcommand{\Set}{{\bf Set}}
\newcommand{\Vect}{{\bf Vect}}
\newcommand{\FDVect}{{\bf FDVect}}
%\newcommand{\Ring}{{\bf Ring}}
\newcommand{\Ab}{{\bf Ab}}
\newcommand{\Mod}{\mathop{\mathrm{Mod}}\nolimits}
\newcommand{\Modf}{\mathop{\mathrm{Mod}^\mathrm{f}}\nolimits}
\newcommand{\CGA}{{\bf CGA}}
\newcommand{\GVect}{{\bf GVect}}
\newcommand{\Lie}{{\bf Lie}}
\newcommand{\dLie}{{\bf Liec}}



% 射の集合など
\newcommand{\Map}{\mathop{\mathrm{Map}}\nolimits} % 写像の集合
\newcommand{\Hom}{\mathop{\mathrm{Hom}}\nolimits} % 射集合
\newcommand{\End}{\mathop{\mathrm{End}}\nolimits} % 自己準同型の集合
\newcommand{\Aut}{\mathop{\mathrm{Aut}}\nolimits} % 自己同型の集合
\newcommand{\Mor}{\mathop{\mathrm{Mor}}\nolimits} % 射集合
\newcommand{\Ker}{\mathop{\mathrm{Ker}}\nolimits} % 核
\newcommand{\Img}{\mathop{\mathrm{Im}}\nolimits} % 像
\newcommand{\Cok}{\mathop{\mathrm{Coker}}\nolimits} % 余核
\newcommand{\Cim}{\mathop{\mathrm{Coim}}\nolimits} % 余像

% その他便利なコマンド
\newcommand{\dip}{\displaystyle} % 本文中で数式モード
\newcommand{\e}{\varepsilon} % イプシロン
\newcommand{\dl}{\delta} % デルタ
\newcommand{\pphi}{\varphi} % ファイ
\newcommand{\ti}{\tilde} % チルダ
\newcommand{\pal}{\parallel} % 平行
\newcommand{\op}{{\rm op}} % 双対を取る記号
\newcommand{\lcm}{\mathop{\mathrm{lcm}}\nolimits} % 最小公倍数の記号
\newcommand{\Probsp}{(\Omega, \F, \P)} 
\newcommand{\argmax}{\mathop{\rm arg~max}\limits}
\newcommand{\argmin}{\mathop{\rm arg~min}\limits}





% ================================
% コマンドを追加する場合のスペース 
%\newcommand{\OO}{\mcal{O}}



\makeatletter
\renewenvironment{proof}[1][\proofname]{\par
  \pushQED{\qed}%
  \normalfont \topsep6\p@\@plus6\p@\relax
  \trivlist
  \item[\hskip\labelsep
%        \itshape
         \bfseries
%    #1\@addpunct{.}]\ignorespaces
    {#1}]\ignorespaces
}{%
  \popQED\endtrivlist\@endpefalse
}
\makeatother

\renewcommand{\proofname}{証明.}



%\renewcommand\proofname{\bf 証明} % 証明
\numberwithin{equation}{section}
\newcommand{\cTop}{\textsf{Top}}
%\newcommand{\cOpen}{\textsf{Open}}
\newcommand{\Op}{\mathop{\textsf{Op}}\nolimits}
\newcommand{\Ob}{\mathop{\textrm{Ob}}\nolimits}
\newcommand{\id}{\mathop{\mathrm{id}}\nolimits}
\newcommand{\pt}{\mathop{\mathrm{pt}}\nolimits}
\newcommand{\res}{\mathop{\rho}\nolimits}
\newcommand{\A}{\mcal{A}}
\newcommand{\B}{\mcal{B}}
\newcommand{\C}{\mcal{C}}
\newcommand{\D}{\mcal{D}}
\newcommand{\E}{\mcal{E}}
\newcommand{\G}{\mcal{G}}
%\newcommand{\H}{\mcal{H}}
\newcommand{\I}{\mcal{I}}
\newcommand{\J}{\mcal{J}}
\newcommand{\OO}{\mcal{O}}
\newcommand{\Ring}{\mathop{\textsf{Ring}}\nolimits}
\newcommand{\cAb}{\mathop{\textsf{Ab}}\nolimits}
%\newcommand{\Ker}{\mathop{\mathrm{Ker}}\nolimits}
\newcommand{\im}{\mathop{\mathrm{Im}}\nolimits}
\newcommand{\Coker}{\mathop{\mathrm{Coker}}\nolimits}
\newcommand{\Coim}{\mathop{\mathrm{Coim}}\nolimits}
\newcommand{\rank}{\mathop{\mathrm{rank}}\nolimits}
\newcommand{\Ht}{\mathop{\mathrm{Ht}}\nolimits}
\newcommand{\supp}{\mathop{\mathrm{supp}}\nolimits}
\newcommand{\colim}{\mathop{\mathrm{colim}}}
\newcommand{\Tor}{\mathop{\mathrm{Tor}}\nolimits}

\newcommand{\cat}{\mathscr{C}}

%筆記体
\newcommand{\cA}{\mcal{A}}
\newcommand{\cB}{\mcal{B}}
\newcommand{\cC}{\mcal{C}}
\newcommand{\cD}{\mcal{D}}
\newcommand{\cE}{\mcal{E}}
\newcommand{\cF}{\mcal{F}}
\newcommand{\cG}{\mcal{G}}
\newcommand{\cH}{\mcal{H}}
\newcommand{\cI}{\mcal{I}}
\newcommand{\cJ}{\mcal{J}}
\newcommand{\cK}{\mcal{K}}
\newcommand{\cL}{\mcal{L}}
\newcommand{\cM}{\mcal{M}}
\newcommand{\cN}{\mcal{N}}
\newcommand{\cO}{\mcal{O}}
\newcommand{\cP}{\mcal{P}}
\newcommand{\cQ}{\mcal{Q}}
\newcommand{\cR}{\mcal{R}}
\newcommand{\cS}{\mcal{S}}
\newcommand{\cT}{\mcal{T}}
\newcommand{\cU}{\mcal{U}}
\newcommand{\cV}{\mcal{V}}
\newcommand{\cW}{\mcal{W}}
\newcommand{\cX}{\mcal{X}}
\newcommand{\cY}{\mcal{Y}}
\newcommand{\cZ}{\mcal{Z}}


\newcommand{\scA}{\mathscr{A}}
\newcommand{\scB}{\mathscr{B}}
\newcommand{\scC}{\mathscr{C}}
\newcommand{\scD}{\mathscr{D}}
\newcommand{\scE}{\mathscr{E}}
\newcommand{\scF}{\mathscr{F}}
\newcommand{\scN}{\mathscr{N}}
\newcommand{\scO}{\mathscr{O}}
\newcommand{\scV}{\mathscr{V}}
\newcommand{\scU}{\mathscr{U}}


\newcommand{\ibA}{\mathop{\text{\textit{\textbf{A}}}}}
\newcommand{\ibB}{\mathop{\text{\textit{\textbf{B}}}}}
\newcommand{\ibC}{\mathop{\text{\textit{\textbf{C}}}}}
\newcommand{\ibD}{\mathop{\text{\textit{\textbf{D}}}}}
\newcommand{\ibE}{\mathop{\text{\textit{\textbf{E}}}}}
\newcommand{\ibF}{\mathop{\text{\textit{\textbf{F}}}}}
\newcommand{\ibG}{\mathop{\text{\textit{\textbf{G}}}}}
\newcommand{\ibH}{\mathop{\text{\textit{\textbf{H}}}}}
\newcommand{\ibI}{\mathop{\text{\textit{\textbf{I}}}}}
\newcommand{\ibJ}{\mathop{\text{\textit{\textbf{J}}}}}
\newcommand{\ibK}{\mathop{\text{\textit{\textbf{K}}}}}
\newcommand{\ibL}{\mathop{\text{\textit{\textbf{L}}}}}
\newcommand{\ibM}{\mathop{\text{\textit{\textbf{M}}}}}
\newcommand{\ibN}{\mathop{\text{\textit{\textbf{N}}}}}
\newcommand{\ibO}{\mathop{\text{\textit{\textbf{O}}}}}
\newcommand{\ibP}{\mathop{\text{\textit{\textbf{P}}}}}
\newcommand{\ibQ}{\mathop{\text{\textit{\textbf{Q}}}}}
\newcommand{\ibR}{\mathop{\text{\textit{\textbf{R}}}}}
\newcommand{\ibS}{\mathop{\text{\textit{\textbf{S}}}}}
\newcommand{\ibT}{\mathop{\text{\textit{\textbf{T}}}}}
\newcommand{\ibU}{\mathop{\text{\textit{\textbf{U}}}}}
\newcommand{\ibV}{\mathop{\text{\textit{\textbf{V}}}}}
\newcommand{\ibW}{\mathop{\text{\textit{\textbf{W}}}}}
\newcommand{\ibX}{\mathop{\text{\textit{\textbf{X}}}}}
\newcommand{\ibY}{\mathop{\text{\textit{\textbf{Y}}}}}
\newcommand{\ibZ}{\mathop{\text{\textit{\textbf{Z}}}}}

\newcommand{\ibx}{\mathop{\text{\textit{\textbf{x}}}}}

%\newcommand{\Comp}{\mathop{\mathrm{C}}\nolimits}
%\newcommand{\Komp}{\mathop{\mathrm{K}}\nolimits}
%\newcommand{\Domp}{\mathop{\mathsf{D}}\nolimits}%複体のホモトピー圏
%\newcommand{\Comp}{\mathrm{C}}
%\newcommand{\Komp}{\mathrm{K}}
%\newcommand{\Domp}{\mathsf{D}}%複体のホモトピー圏

\newcommand{\Comp}{\mathop{\mathrm{C}}\nolimits}
\newcommand{\Komp}{\mathop{\mathsf{K}}\nolimits}
\newcommand{\Domp}{\mathord{\mathsf{D}}}%\nolimits}
\newcommand{\Kompl}{\mathop{\mathsf{K}^\mathrm{+}}\nolimits}
\newcommand{\Kompu}{\mathop{\mathsf{K}^\mathrm{-}}\nolimits}
\newcommand{\Kompb}{\mathop{\mathsf{K}^\mathrm{b}}\nolimits}
\newcommand{\Dompl}{\mathop{\mathsf{D}^\mathrm{+}}\nolimits}
\newcommand{\Dompu}{\mathop{\mathsf{D}^\mathrm{-}}\nolimits}
\newcommand{\Dompb}{\mathop{\mathsf{D}^\mathrm{b}}\nolimits}
\newcommand{\Dbconic}{\mathop{\mathsf{D}^\mathrm{b}_{\rrp}}\nolimits}
\newcommand{\Dompbf}{\mathop{\mathsf{D}_\mathrm{f}^\mathrm{b}}\nolimits}
\newcommand{\Domplb}{\mathop{\mathsf{D}^\mathrm{lb}}\nolimits}




\newcommand{\CCat}{\Comp(\cat)}
\newcommand{\KCat}{\Komp(\cat)}
\newcommand{\DCat}{\Domp(\cat)}%圏Cの複体のホモトピー圏
\newcommand{\HOM}{\mathop{\mathscr{H}\hspace{-2pt}om}\nolimits}%内部Hom
\newcommand{\RHOM}{\mathop{\mathrm{R}\hspace{-1.5pt}\HOM}\nolimits}

\newcommand{\muS}{\mathop{\mathrm{SS}}\nolimits}
\newcommand{\RG}{\mathop{\mathrm{R}\hspace{-0pt}\Gamma}\nolimits}
\newcommand{\RHom}{\mathop{\mathrm{R}\hspace{-1.5pt}\Hom}\nolimits}
\newcommand{\Rder}{\mathrm{R}}

\newcommand{\simar}{\mathrel{\overset{\sim}{\rightarrow}}}%同型右矢印
\newcommand{\simarr}{\mathrel{\overset{\sim}{\longrightarrow}}}%同型右矢印
\newcommand{\simra}{\mathrel{\overset{\sim}{\leftarrow}}}%同型左矢印
\newcommand{\simrra}{\mathrel{\overset{\sim}{\longleftarrow}}}%同型左矢印

\newcommand{\hocolim}{{\mathrm{hocolim}}}
\newcommand{\indlim}[1][]{\mathop{\varinjlim}\limits_{#1}}
\newcommand{\sindlim}[1][]{\smash{\mathop{\varinjlim}\limits_{#1}}\,}
\newcommand{\Pro}{\mathrm{Pro}}
\newcommand{\Ind}{\mathrm{Ind}}
\newcommand{\prolim}[1][]{\mathop{\varprojlim}\limits_{#1}}
\newcommand{\sprolim}[1][]{\smash{\mathop{\varprojlim}\limits_{#1}}\,}

\newcommand{\Sh}{\mathrm{Sh}}
\newcommand{\PSh}{\mathrm{PSh}}

\newcommand{\rmD}{\mathrm{D}}

\newcommand{\Lloc}[1][]{\mathord{\mathcal{L}^1_{\mathrm{loc},{#1}}}}
\newcommand{\ori}{\mathord{\mathrm{or}}}
\newcommand{\Db}{\mathord{\mathscr{D}b}}

\newcommand{\codim}{\mathop{\mathrm{codim}}\nolimits}



\newcommand{\gld}{\mathop{\mathrm{gld}}\nolimits}
\newcommand{\wgld}{\mathop{\mathrm{wgld}}\nolimits}


\newcommand{\tens}[1][]{\mathbin{\otimes_{\raise1.5ex\hbox to-.1em{}{#1}}}}
\newcommand{\etens}{\mathbin{\boxtimes}}
\newcommand{\ltens}[1][]{\mathbin{\overset{\mathrm{L}}\tens}_{#1}}
\newcommand{\mtens}[1][]{\mathbin{\overset{\mathrm{\mu}}\tens}_{#1}}
\newcommand{\lltens}[1][]{{\mathop{\tens}\limits^{\mathrm{L}}_{#1}}}
\newcommand{\letens}{\overset{\mathrm{L}}{\etens}}
\newcommand{\detens}{\underline{\etens}}
\newcommand{\ldetens}{\overset{\mathrm{L}}{\underline{\etens}}}
\newcommand{\dtens}[1][]{{\overset{\mathrm{L}}{\underline{\otimes}}}_{#1}}

\newcommand{\blk}{\mathord{\ \cdot\ }}
\newcommand{\mres}[2][]{{\left.{#1}\right\rvert}_{#2}}


%\newcommand{\hocolim}{{\mathrm{hocolim}}}
%\newcommand{\indlim}[1][]{\mathop{\varinjlim}\limits_{#1}}
%\newcommand{\sindlim}[1][]{\smash{\mathop{\varinjlim}\limits_{#1}}\,}
%\newcommand{\Pro}{\mathrm{Pro}}
%\newcommand{\Ind}{\mathrm{Ind}}
%\newcommand{\prolim}[1][]{\mathop{\varprojlim}\limits_{#1}}
%\newcommand{\sprolim}[1][]{\smash{\mathop{\varprojlim}\limits_{#1}}\,}
\newcommand{\proolim}[1][]{\mathop{\text{\rm``{$\varprojlim$}''}}\limits_{#1}}
\newcommand{\sproolim}[1][]{\smash{\mathop{\rm``{\varprojlim}''}\limits_{#1}}}
\newcommand{\inddlim}[1][]{\mathop{\text{\rm``{$\varinjlim$}''}}\limits_{#1}}
\newcommand{\sinddlim}[1][]{\smash{\mathop{\text{\rm``{$\varinjlim$}''}}\limits_{#1}}\,}
\newcommand{\ooplus}{\mathop{\text{\rm``{$\oplus$}''}}\limits}
\newcommand{\bbigsqcup}{\mathop{``\bigsqcup''}\limits}
\newcommand{\bsqcup}{\mathop{``\sqcup''}\limits}
\newcommand{\dsum}[1][]{\mathbin{\oplus_{#1}}}

\newcommand{\Fct}{\mathop{\mathsf{Fct}}\nolimits}

\newcommand{\pb}[1][]{\mathbin{\mathop{\times}\limits_{#1}}}





%================================================
% 自前の定理環境
%   https://mathlandscape.com/latex-amsthm/
% を参考にした
\newtheoremstyle{mystyle}%   % スタイル名
    {5pt}%                   % 上部スペース
    {5pt}%                   % 下部スペース
    {}%              % 本文フォント
    {}%                  % 1行目のインデント量
    {\bfseries}%                      % 見出しフォント
    {.}%                     % 見出し後の句読点
    {12pt}%                     % 見出し後のスペース
    {\thmname{#1}\thmnumber{ #2}\thmnote{{\hspace{2pt}\normalfont (#3)}}}% % 見出しの書式

\theoremstyle{mystyle}
\newtheorem{AXM}{公理}[section]
\newtheorem{DFN}[AXM]{定義}
\newtheorem{THM}[AXM]{定理}
\newtheorem*{THM*}{定理}
\newtheorem{PRP}[AXM]{命題}
\newtheorem{LMM}[AXM]{補題}
\newtheorem{CRL}[AXM]{系}
\newtheorem{EG}[AXM]{例}
\newtheorem*{EG*}{例}
\newtheorem{RMK}[AXM]{注意}
\newtheorem{CNV}[AXM]{約束}
\newtheorem{CMT}[AXM]{コメント}
\newtheorem*{CMT*}{コメント}
\newtheorem{NTN}[AXM]{記号}

% 定理環境ここまで
%====================================================

\usepackage{framed}
\definecolor{lightgray}{rgb}{0.75,0.75,0.75}
\renewenvironment{leftbar}{%
  \def\FrameCommand{\textcolor{lightgray}{\vrule width 4pt} \hspace{10pt}}% 
  \MakeFramed {\advance\hsize-\width \FrameRestore}}%
{\endMakeFramed}
\newenvironment{redleftbar}{%
  \def\FrameCommand{\textcolor{lightgray}{\vrule width 1pt} \hspace{10pt}}% 
  \MakeFramed {\advance\hsize-\width \FrameRestore}}%
 {\endMakeFramed}



\renewcommand{\Re}{\mathop{\mathrm{Re}}\nolimits}

% =================================





% ---------------------------
% new definition macro
% ---------------------------
% 便利なマクロ集です

% 内積のマクロ
%   例えば \inner<\pphi | \psi> のように用いる
\def\inner<#1>{\langle #1 \rangle}

% ================================
% マクロを追加する場合のスペース 

%=================================



% 位相空間
\newcommand{\Int}[1][]{\mathop{\mathrm{Int}}\nolimits_{#1}}
\newcommand{\Cl}[1][]{\mathop{\mathrm{Cl}}\nolimits_{#1}}
\newcommand{\Bd}[1][]{\mathop{\mathrm{Bd}}\nolimits_{#1}}
\newcommand{\bd}{\mathord{\partial}}
\newcommand{\Fr}[1][]{\mathop{\mathrm{Fr}}\nolimits_{#1}}


\newcommand{\mtan}[1][]{T\hspace{-1.75pt}{#1}}
\newcommand{\mtp}[1][]{{}^{t\!}{#1}} % transpose of the morphism
%\newcommand{\mpb}[1][]{{}^{t\!}{#1}'} % pull-back of the morphism
\newcommand{\mpb}[1][]{{#1}_{\pi}} % pull-back of the morphism
\newcommand{\mdiff}[1][]{{#1}_{d}} % pull-back of the morphism
\newcommand{\mptan}[1][]{{#1}_{\tau}} % pull-back of the morphism


%\newcommand{\pb}[1][]{\mathbin{\mathop{\times}\limits_{#1}}}

\newcommand{\cotb}[1][]{T^{\ast}{#1}}
\newcommand{\conm}[2][]{T_{#1}^{\hspace{0.7pt}\ast}{#2}}
\newcommand{\norb}[2][]{T_{#1}{#2}}

\newcommand{\cotz}[1][]{T^{\ast}_{#1}{#1}}

\newcommand{\cotn}[1][]{\dot{T}^{\ast}{#1}}


\newcommand{\muhom}{\mu hom}
\newcommand{\muHom}[1][]{\mathrm{Hom}^\mu_{\raise1.5ex\hbox to.1em{}#1}}

\newcommand{\mucat}[2][]{\mathsf{D}^{\mathrm{b}}_{#1}(#2)}
%\newcommand{\mucat}[2][]{\mathsf{D}^{\mathrm{b}}_{#1}(#2)}

\newcommand{\conv}[1][]{\mathbin{\mathop{\circ}\limits_{#1}}}

\newcommand{\torus}{\mathbf{T}}



% ----------------------------
% documenet 
% ----------------------------
% 以下, 本文の執筆スペースです. 
% Your main code must be written between 
% begin document and end document.
% ---------------------------

\title{ゼミノート（4月16日分）}
\author{Toshi2019}
\date{\today}
\begin{document}
\maketitle
\tableofcontents

\section{（余）接束の基本図式}
ここでは多様体は\(C^\infty\)級または\(C^\omega\)級とする．

\(f\colon Y\to X\)を多様体の射とする．
\(\tau\colon TX\to X\)等を接束とすると，
ベクトル束の引き戻し
\[\xymatrix{
  Y\pb[X]TX
  \ar[r]^-{f_{\tau}}
  \ar[d]^-{\tau}
  \ar@{}[dr]|\square
  &
  TX
  \ar[d]^-{\tau}
  \\
  Y\ar[r]^-{f}&X
}\]
が定まる．\(f\)の微分を\(df\colon TY\to TX\)で表すとき，
図式
\[\xymatrix{
  TY
  \ar@/^1pc/[rrd]^-{df}
  \ar@/_1pc/[rdd]_-{\tau}\\
  &Y\pb[X]TX
  \ar[r]^-{f_{\tau}}
  \ar[d]^-{\tau}
  \ar@{}[dr]|\square
  &
  TX
  \ar[d]^-{\tau}
  \\
  &Y\ar[r]^-{f}&X
}\]
は可換である．
したがって，引き戻しの普遍性から
射\(f'\colon TY\to Y\mathbin{\times_{X}}TX\)がひきおこされ，
\(df\)は\(df=f_{\tau}\circ f'\)と分解される．
\begin{equation}\vcenter{
\xymatrix{
  TY
  \ar@/^2pc/[rr]^-{df}
  \ar[r]^-{f'}
  \ar[rd]_-{\tau}
  &Y\pb[X]TX
  \ar[r]^-{f_{\tau}}
  \ar[d]^-{\tau}
  \ar@{}[dr]|\square
  &
  TX
  \ar[d]^-{\tau}
  \\
  &Y\ar[r]^-{f}&X
}}\end{equation}
この図式を，ここだけの用語で，接束の基本図式とよぶことにする．

\(\pi\colon T^{\ast}X\to X\)等を余接束とする．
先と同様にベクトル束の引き戻し
\[\xymatrix{
  Y\pb[X]T^{\ast}X
  \ar[r]^-{f_{\pi}}
  \ar[d]^-{\pi}
  \ar@{}[dr]|\square
  &
  T^{\ast}X
  \ar[d]^-{\pi}
  \\
  Y\ar[r]^-{f}&X
}\]
が定まる．
他方で，\(Y\)上の束の射\(
  f'\colon TY\to Y\mathbin{\times_{X}}TX
\)の双対射\(
  {}^{t\!}f'\colon \left(Y\mathbin{\times_{X}}TX\right)^{\ast}\to T^{\ast}Y
\)が定まる．
\(\left(Y\mathbin{\times_{X}}TX\right)^{\ast}\)
は
\(Y\mathbin{\times_{X}}T^{\ast}X\)と同型なので，
射\(
  f_{d}\colon Y\mathbin{\times_{X}}T^{\ast}X\to T^{\ast}Y
\)が定まる．
これにより，図式
\begin{equation}\label{eq:cot}
\vcenter{\xymatrix{
  T^{\ast}Y
  %\ar@/^2pc/[rr]^-{df}
  \ar[rd]_-{\pi}
  &
  Y\pb[X]T^{\ast}X
  \ar[l]_-{f_{d}}
  \ar[r]^-{f_{\pi}}
  \ar[d]^-{\pi}
  \ar@{}[dr]|\square
  &
  T^{\ast}X
  \ar[d]^-{\pi}
  \\
  &Y\ar[r]^-{f}&X
}}\end{equation}
を得る．
この図式を，ここだけの用語で，余接束の基本図式とよぶことにする．

\section{余接束のシンプレクティック構造}
\subsection*{標準1形式}

\eqref{eq:cot}を\(\pi\colon \cotb[X]\to X\)に適用する．
\begin{equation}\label{eq:cotcot}
    \vcenter{\xymatrix{
      T^{\ast}\cotb[X]
      %\ar@/^2pc/[rr]^-{df}
      \ar[rd]_-{\pi}
      &
      \cotb[X]\pb[X]\cotb[X]
      \ar[l]_-{\pi_{d}}
      \ar[r]^-{\pi_{\pi}}
      \ar[d]^-{\pi}
      \ar@{}[dr]|\square
      &
      \cotb[X]
      \ar[d]^-{\pi}
      \\
      &\cotb[X]\ar[r]^-{\pi}&X
}}\end{equation}
\(\cotb[X]\mathbin{\times_{X}}\cotb[X]\)の対角集合を\(
    \varDelta
\)とする．\(\varDelta\)と\(\cotb[X]\)を同一視した上で，
\(\pi_{d}\)を\(\varDelta\)に制限することで
\(\pi\colon\cotb[{\cotb[X]}]\to\cotb[X]\)の切断
\[
    \mres[\pi_{d}]{\varDelta}\colon
    \cotb[X]
    \cong\varDelta
    \to\cotb[{\cotb[X]}]
\]    
が得られる．
これを\(\cotb[X]\)の\textbf{標準1形式} (canonical 1-form) といい，
\(\alpha_X\)で表す．

\subsection*{標準1形式の座標表示}

\((x;\xi)\)を\(\cotb[X]\)の局所座標とする．
このとき，\(\alpha_X\)は
\[
    \alpha_{X}(x;\xi)=((x;\xi);\pi^{\ast}\xi)
\]
とかける．ここで，\[
    \pi^{\ast}\xi\in T^{\ast}_{(x;\xi)}T^{\ast}X
\]は\(T_{(x;\xi)}T^{\ast}X\)の基底\(
    \frac{\partial}{\partial x_{1}},
    \dots,
    \frac{\partial}{\partial x_{n}},
    \frac{\partial}{\partial \xi_{1}},
    \dots,
    \frac{\partial}{\partial \xi_{n}}
\)に関して，
\begin{align*}
    \left\langle \pi^{\ast}\xi, \frac{\partial}{\partial x_{j}}\right\rangle
    &=
    \left\langle \xi, d\pi_{(x;\xi)}\frac{\partial}{\partial x_{j}}\right\rangle\\
    &=
    \left\langle \xi, \frac{\partial}{\partial x_{j}}\right\rangle,\\
    \left\langle \pi^{\ast}\xi, \frac{\partial}{\partial \xi_{j}}\right\rangle
    &=
    \left\langle \xi, d\pi_{(x;\xi)}\frac{\partial}{\partial \xi_{j}}\right\rangle\\
    &=
    \left\langle \xi, 0\right\rangle=0
\end{align*}
をみたす．したがって，\(\alpha_X\)は基底\(dx_{1},\dots,dx_{n},d\xi_{1},\dots,d\xi_{n}\)を用いて，
\[
    \alpha_X
    =\pi^{\ast}\left(
        \sum_{j=1}^{n}\xi_{j}dx_{j}
    \right)
    =\sum_{j=1}^{n}\xi_{j}\pi^{\ast}dx_{j}
    =\sum_{j=1}^{n}\xi_{j}dx_{j}
\]
と表せる．

\subsection*{余接束上のシンプレクティック形式}
以上で構成した\(\alpha_X\)に対し，\[\omega_X=d\alpha_X\]とおく．
\(\omega_X\)は\(\cotb[X]\)上のシンプレクティック形式である．したがって，
\((\cotb[X],\omega_X)\)はシンプレクティック多様体となる．

\(\omega_X\)を座標表示すると，
\[
    \omega_X=d\alpha_{X}
    =d\left(\sum_{j=1}^{n}\xi_{j}dx_{j}\right)
    =\sum_{j=1}^{n}d(\xi_{j}dx_{j})
    =\sum_{j=1}^{n}d\xi_{j}\wedge dx_{j}
\]
となる．

\section{錐状層の超局所台 (pp.241--245)}
\newcommand{\eu}{\mathbf{e}}
\newcommand{\rrp}{\rr^{+}}

この節ではベクトル束上の層の超局所台について調べる．

\(\tau\colon E\to Z\)を多様体\(Z\)上の
実ベクトル束とする（III \S7参照）．
\(\rrp\)の\(E\)への作用を表す写像を\(
    \mu\colon \rrp\times E\to E
\)で表し，
\(\mu\)の無限小作用を表す\(E\)上のベクトル場
を\(\eu\)で表す．
すなわち\(E\)上の任意の関数\(\varphi\)に対し，\(
    (\eu(\varphi))(x)
    =
    \dfrac{\partial}{\partial{t}}\varphi(\mu(t,x))|_{t=1}
\)が成り立つ．
ベクトル場\(\eu\)は
\textbf{オイラーベクトル場} (the Euler vector field) と
呼ばれることが多い．
\begin{CMT}
    \(\eu\)を作用素とみる．
    \(\varphi\in C^{\infty}(E)\)に対し，\(
        \varphi\circ \mu\in C^{\infty}(\rrp \times E)
    \)としてから，\(t=1\)での偏微分\[
        \ev_{t=1}\left(\frac{\partial}{\partial t}(\varphi\circ \mu)\right)
    \]を対応させることで，\(\eu(\varphi)\)が定まる．
\end{CMT}
写像\(\mu\)から写像\(
    \mdiff[\mu]\colon
    \rrp\times\cotb[E]\to\cotb[\rrp]\times\cotb[E]
\)が定まる．\(\mdiff[\mu]\)の\(1\in\rrp\)への制限を考え，射影\(
    \cotb[\rrp]\times\cotb[E]
    \to
    \cotb[\rrp]\cong\rrp\times\rr\to\rr
\)と合成することにより，写像\begin{equation}
    \theta_{E}\colon\cotb[E]\to\rr
\end{equation}を得る．
この写像は\(\eu\)の主要表象に他ならない．
\begin{CMT}
    \(\mu\colon \rrp\times E\to E\)に関する余接束の基本図式を考える．
    \[\vcenter{\xymatrix{
        \cotb[(\rrp\times E)]
        \ar[rd]_-{\pi}
        &
        (\rrp\times E)\pb[E]\cotb[E]
        \ar[l]_-{\mdiff[\mu]}
        \ar[r]^-{\mpb[\mu]}
        \ar[d]^-{\pi}
        &
        \cotb[E]
        \ar[d]^-{\pi}
        \\
        &
        \rrp\times E
        \ar[r]^-{\mu}
        &
        E
    }}\]
    \(\rrp\times \cotb[E]\)から
    \(\rrp\times E\)と\(\cotb[E]\)への射を
    \begin{align*}
        \id_{\rrp}\times\pi&\colon
        \rrp\times \cotb[E] \to\rrp\times E,\\
        \widetilde{\mu}&\colon
        \rrp\times \cotb[E] \to \cotb[E];
        \quad 
        (t,x;\mu^{\ast}\xi)\mapsto(\mu(t,x);\xi)
    \end{align*}
    で定めると，\(\rrp\times \cotb[E]\)も引き戻しの普遍性をみたすので，
    \[
        (\rrp\times E)\pb[E]\cotb[E]
        \cong
        \rrp\times \cotb[E]
    \]が成り立つ．
    この同型は，\[
        (\id_{\rrp}\times\pi,\widetilde{\mu})\colon
        \rrp\times \cotb[E]
        \to
        (\rrp\times E)\pb[E]\cotb[E];\quad
        \left(t,\left(x;\xi\right)\right)
        \mapsto
        ((t,x),(\mu(t,x);\xi))
    \]と\[
        (\id_{\rrp}\times\pi,\widetilde{\mu})\colon
        (\rrp\times E)\pb[E]\cotb[E]
        \to
        \rrp\times \cotb[E]
        ;\quad
        \left(t,\left(x;\xi\right)\right)
        \mapsto
        ((t,x),(\mu(t,x);\xi))
    \]とで与えられる．
    
    \(
        \mdiff[\mu]\colon
        \rrp\times \cotb[E]
        \cong
        (\rrp\times E)\times_{E}\cotb[E]
        \to
        \cotb[\rrp]\times\cotb[E]
    \)の\(t=1\in\rrp\)への制限\(
        \mres[{\mdiff[\mu]}]{\{t=1\}}
        \colon
        \{1\}\times \cotb[E]\to
        \cotb[\rrp]\times\cotb[E]
    \)を考える．
\end{CMT}

\((z,x)\)を（\(Z\)上局所的な）座標系で，
\((z)\)が\(Z\)の座標で\((x)\)が線形座標であるものとする．
\((z,x;\zeta,\xi)\)で対応する\(\cotb[E]\)の座標を表す．
このとき，
\begin{equation}
    \begin{dcases}
        \eu=\left\langle{x,\dfrac{\partial}{\partial{x}}}\right\rangle=\sum_{j}x_{j}\dfrac{\partial}{\partial{x_{j}}},\\
        \theta_{E}=\left\langle{x,\xi}\right\rangle,\\
        \alpha_{E}=\left\langle{\zeta,dz}\right\rangle+\left\langle{\xi,dx}\right\rangle
    \end{dcases}
\end{equation}
である．
ただし，\(\alpha_{E}\)は\(\cotb[E]\)の標準\(1\)形式である
（A\S2参照）．

\(\cotb[(E/Z)]\)で\(Z\)上の\(E\)の相対余接束を表していた．
これは\(E\)上のベクトル束の完全列
\begin{equation}
    0\to E\pb[Z]\cotb[Z]\to\cotb[E]\to\cotb[(E/Z)]\to0
\end{equation}
で定義される．
任意の\(z\in Z\)に対し，
\(\cotb[(E/Z)]\)の\(z\)上のファイバー\(\cotb[(E/Z)]_{z}\)は
\(E\)の\(z\)上のファイバー\(E_{z}\)の余接束\(\cotb[(E_{z})]\)と
同一視できる．
他方で，\(\cotb[(E_{z})]\)は\(E_{z}\times E^{\ast}_{z}\)と
同型である．
したがって，各\(z\)に対し，
\(\cotb[(E/Z)]_{z}\)から\(E^{\ast}_{z}\)への写像が定まる．
これから写像\begin{equation}
    \cotb[(E/Z)]\to E^{\ast}
\end{equation}がえられる．このとき\(Z\)上の射
\begin{equation}
    \psi_{E}\colon\cotb[E]\to E^{\ast}
\end{equation}
を\(\cotb[E]\to \cotb[(E/Z)]\to E^{\ast}\)の合成として定める．

\begin{leftbar}
\begin{PRP}\label{prp551}
    \begin{enumerate}[(i)]
        \item 写像\(\varPhi\colon \cotb[E]\to\cotb[E]\)で，
        \(
            \alpha_{E}-d\theta_{E}
            =\varPhi_{E}^{\ast}(\alpha_{E^{\ast}})
        \)であり\(
            \cotb[E]\underset{\varPhi_{E}}{\rightarrow}
            \cotb[E^{\ast}]\underset{\alpha}{\rightarrow}
            E^{\ast}
        \)の合成が\(\psi_{E}\)となるものがただひとつ存在する．
        \item \(\theta_{E^{\ast}}\circ\varPhi_{E}=-\theta_E\)である．
        \item \(\varPhi_{E^{\ast}}\circ \varPhi_{E}=-a^{\ast}\)である．ただし，\(a^{\ast}\)は\(E\)上の対蹠写像からひきおこされる\(\cotb[E]\)の自己同型である．
    \end{enumerate}
\end{PRP}
\end{leftbar}

\begin{proof}
    \begin{equation}
        \varPhi_{E}\colon (z,x;\zeta,\xi)\mapsto(z,\xi;\zeta,-x)
    \end{equation}
\end{proof}

\begin{RMK}\label{rmk552}
    \(\varPhi_{E}\)はシンプレクティック変換だが
    斉次シンプレクティック変換ではない．
\end{RMK}

\begin{equation}
    S_{E}=\theta_{E}^{-1}(0)
\end{equation}

\begin{equation}
    \begin{dcases}
        (z,x;\zeta,\xi)\mapsto (z,x;\lambda\zeta,\lambda\xi),& \lambda\in\rrp,\\
        (z,x;\zeta,\xi)\mapsto (z,\mu x;\zeta,\mu^{-1}\xi),& \mu\in\rrp.
    \end{dcases}
\end{equation}

\begin{leftbar}
\begin{PRP}\label{prp553}
    \begin{enumerate}[(i)]
        \item \(\Dbconic(E)\)
        \item \(F\in\Dbconic(E)\)のとき，\(\muS(F)\)は2重錐状
    \end{enumerate}
\end{PRP}
\end{leftbar}
    
\begin{equation}
    \vcenter{\xymatrix{
        \cotb[E]
        &
        E\pb[Z]\cotb[Z]
        \ar@{_{(}->}[l]-^{\mdiff[\tau]}
        &
        \cotb[Z]\\
        \ar@{_{(}->}[u]-^{\mdiff[\dot{\tau}]}
        \cotb[{\dot{E}}]
        &
        \ar@{_{(}->}[u]-^{\mdiff[\dot{\tau}]}
        \ar@{_{(}->}[l]-^{\mdiff[\dot{\tau}]}
        \dot{E}\pb[Z]\cotb[Z]
        &
        \ar@{_{(}->}[u]-^{\mdiff[\dot{\tau}]}
        \cotb[Z].
    }}
\end{equation}

\begin{equation}
    \cotb[Z]\hookrightarrow
    E\pb[Z]\cotb[Z]
    \underset{\mdiff[\tau]}{\hookrightarrow}
    \cotb[E].
\end{equation}



\begin{equation}
    \mpb[\tau]\mdiff[\tau]^{-1}(A)=\cotb[Z]\cap A.
\end{equation}

\begin{leftbar}
\begin{PRP}
    \(F\in\Dbconic(E)\)とする．このとき，次がなりたつ．
    \begin{enumerate}[(i)]
        \item \(\muS(\Rder{\tau_{\ast}}(F))\subset \cotb[Z]\cap\muS(F)\),
        \item \(\muS(\Rder{\tau_{!}}(F))\subset \cotb[Z]\cap\muS(F)\),
        \item \(\muS(\Rder{\dot{\tau}_{\ast}}(F))\subset \mpb[\dot{\tau}]\mdiff[\dot{\tau}]^{-1}(\muS(F))\),
        \item \(\muS(\Rder{\dot{\tau}_{!}}(F))\subset \mpb[\dot{\tau}]\mdiff[\dot{\tau}]^{-1}(\muS(F))\).
    \end{enumerate}
\end{PRP}
\end{leftbar}


\begin{leftbar}
\begin{THM}
    \(F\in\Dbconic(E)\)とする．
    \(\cotb[E]\)と\(\cotb[E^{\ast}]\)を\(\varPhi_{E}\)で
    同一視することにより，
    次がなりたつ．
    \[
        \muS(F)=\muS(F^{\wedge}).
    \]
\end{THM}
\end{leftbar}

\begin{proof}
    \begin{align}
        x_{\circ}\ne0,
        \quad
        \xi_{\circ}\ne0,\label{eq:5512}\\
        \RG_{Z}(F)=0.\label{eq:5513}    
    \end{align}
\end{proof}























%===============================================
% 参考文献スペース
%===============================================
\begin{thebibliography}{20} 
  \bibitem[dR]{dR} 
  de Rham, \textit{Vari\'et\'es diff\'erentiables}, Hermann, 1953.
  和訳：ド・ラーム, 微分多様体, 東京図書, 1974.
  \bibitem[GKS12]{GKS12} St\'ephane Guillermou, 
  Masaki Kashiwara, Pierre Schapira, 
  \textit{Sheaf Quantization of Hamiltonian Isotopies and Applications to Nondisplaceability Problems}, 
  Duke. Math. J. 161, No. 2, pp.201--245, 2012.  
  \bibitem[GP74]{GP74} Victor Guillemin, Alan Pollack, 
    \textit{Differential Topology}, 
    Prentice-Hall, 1974.
  \bibitem[KS90]{KS90} Masaki Kashiwara, Pierre Schapira, 
    \textit{Sheaves on Manifolds}, 
    Grundlehren der Mathematischen Wissenschaften, 292, Springer, 1990.
  \bibitem[Hat89]{Hat89} 服部晶夫, 多様体 増補版, 岩波全書 288, 岩波書店, 1989.
  %\bibitem[Og02]{Og02} 小木曽啓示, 代数曲線論, 朝倉書店, 2022.
  \bibitem[Hor63]{Hor63} Lars H\"ormander, 
  \textit{Linear Pertial Differential Operators}, Fourth Printing, 
  Grundlehren der Mathematischen Wissenschaften, 116, Springer, 1963.
  \bibitem[Hor71]{Hor71} Lars H\"ormander, 
  \textit{Fourier Integral Operators I}, 
  Acta Math. 127, 79--183, (1971).
  \bibitem[Hor89]{Hor89} Lars H\"ormander, 
  \textit{The Analysis of Linear Pertial Differential Operators I}, Second Edition,
  Grundlehren der Mathematischen Wissenschaften, 256, Springer, 1989.
  \bibitem[Hor85]{Hor85} Lars H\"ormander, 
  \textit{The Analysis of Linear Pertial Differential Operators III}, 
  Grundlehren der Mathematischen Wissenschaften, 274, Springer, 1985.
  \bibitem[Sch85]{Sch85} Pierre Schapira, 
    \textit{Microdifferential Systems in the Complex Domain}, 
    Grundlehren der Mathematischen Wissenschaften, 269, 
    Springer, 1985.
    \bibitem[ST92]{ST92} Pierre Schapira, Nobuyuki Tose, 
    \textit{Morse Inequalities for \(\rr\)-constructible Sheaves}, 
    Adv. Math. 93, pp.1--8, 1992. 
    \bibitem[Zwo12]{Zwo12} Maciej Zworski,
    \textit{Semiclsaaical Analysis}, 
    Graduate Studies in Mathematics, 138, 
    American Mathematical Society, 
    2012. 
\end{thebibliography}

%===============================================




\end{document}
